% Full proof of the isolated nonconstant curve-base candidate.
\begin{proof}[Proof of \cref{prop:p081-curvebase-exact-extension}]
\emph{The completed model is normal and integral.}
The morphism \(\pi\) is flat: locally on the regular base curve its
nonzero coordinate functions act without torsion on the local domains
of the dominant integral \(S\); a torsion-free module over a DVR is
flat \cite[Tag 0539]{Stacks}. Thus \(\mathscr S\) is flat and projective
over \(R\). The local ring \(R_0\) is excellent, and the completion
\(R_0\to R\) is regular
\cite[Tags 07QS and 07GH]{Stacks}. On a finite-type affine chart of
\(S_{R_0}\), its base change to \(R\) is a regular map by
\cite[Tag 07C1]{Stacks}. Both rings are Noetherian; its fibres are
regular, hence normal, so normality ascends by
\cite[Tag 0C22]{Stacks}. This proves normality of \(\mathscr S\).
Here the normality of \(S_{R_0}\) follows by localization from \(S\).
Finally, \(X=Y_K\) is integral by geometric integrality. Any nonempty
affine chart of \(\mathscr S\) injects into its localization at \(t\)
by flatness. This localization is a nonzero domain, since it is an
affine open in the integral \(X\). The chart is therefore a domain.
All these generic opens meet, proving that \(\mathscr S\) is integral.

\emph{The unit ball, not an equality of norm functions.}
For every \(m\geq1\) we claim
\begin{equation}\label{eq:p081-curvebase-unitball}
 H^0(\mathscr S,\mathscr A^{\otimes m})
   =\{s\in H^0(X,L^{\otimes m}):\sup_{\mathcal F}|s|_0\leq1\}.
\end{equation}
Integral sections have norm at most one: take the centre of an analytic
point on the proper \(R\)-model and a local frame there. The section
coefficient belongs to the valuation ring. The model metric makes the
frame have norm one \cite[\S5.3]{BE2021}. This also covers points
whose support is a proper subvariety of \(X\).

Conversely, regard \(s\) as a rational section on the normal integral
\(\mathscr S\). It is regular at every horizontal codimension-one point,
since those points lie in \(X\). At a vertical codimension-one point
\(\zeta_E\), the local ring is a DVR. Put
\(b_E=\operatorname{ord}_E(t)>0\).
The real valuation \(\operatorname{ord}_E/b_E\) on the function field
defines an analytic point \(x_E\in\mathcal F\): it restricts to
\(\operatorname{ord}_t\) on \(K\), since every element of \(R^\times\)
is a unit in that DVR. Completing its valued function field realizes
it as a Berkovich point. In a local frame for \(\mathscr A^{\otimes m}\),
\begin{equation}\label{eq:p081-curvebase-divisorial-value}
             |s|_0(x_E)
                =\mathrm e^{-\operatorname{ord}_E(s)/b_E}.
\end{equation}
This is the normal-model point construction of
\cite[Example~4.11]{BE2021}, here over a discrete valuation ring.
The sup-norm bound implies \(\operatorname{ord}_E(s)\geq0\).
On any affine trivialization, the coefficient of \(s\) therefore lies
in every height-one localization. The normal-domain intersection
criterion \cite[Tag 031T]{Stacks} makes it regular. These local
sections glue, proving \eqref{eq:p081-curvebase-unitball}.

No reducedness of the special fibre was assumed. In particular, this
argument does not claim that the sup norm equals the discrete lattice
norm at all radii; \cite[Theorem~4.19]{BE2021} distinguishes that stronger
assertion from integral closure. Multiplicities \(b_E\) were retained
in \eqref{eq:p081-curvebase-divisorial-value}.

\emph{Integral coefficients in a fixed global monomial system.}
Flat base change, applied only to \(R_0\to R\), gives
\begin{equation}\label{eq:p081-curvebase-basechange}
 H^0(\mathscr S,\mathscr A^{\otimes m})
       =(E_m)_0\otimes_{R_0}R .
\end{equation}
Indeed, this is \cite[Tag 02KH]{Stacks}, with \(\pi\) projective
(in particular quasi-compact and quasi-separated) and the line bundle
quasi-coherent. It is not a claim about base change to the special fibre.
The sections \(q_\ell\) generate \(E_1\). Taking symmetric powers,
using \eqref{eq:p081-curvebase-generation}, and then tensoring with \(R\)
gives a surjection from the finite free \(R\)-module indexed by
degree-\(m\) monomials to the module in
\eqref{eq:p081-curvebase-basechange}. Thus every integral section has
an expression
\begin{equation}\label{eq:p081-curvebase-coefficients}
          a=\sum_{|\nu|=m}c_\nu(t)q^\nu,\qquad c_\nu(t)\in R .
\end{equation}
No uniqueness of these coefficients is needed. As \(q_\ell\) generate
\(A\), their norms are at most one and their maximum is one, both for
the trivial metric on \(S^{\mathrm{an}}\) and for the model metric on
\(\mathcal F\). The same statements hold for all degree-\(m\) monomials.
The two references agree under restriction: a local frame near the
centre on \(S\) pulls back to a local model frame on \(\mathscr S\).

\emph{Truncate the coefficients with a strict pointwise margin.}
Write the weights of \eqref{eq:p081-curvebase-fs} as
\(\lambda_j=n_j+\alpha_j\), where \(n_j\in\mathbf Z\) and
\(0\leq\alpha_j<1\), and set \(a_j=t^{-n_j}s_j\). Then
\[
                \mathrm e^{mf}=\max_j\mathrm e^{\alpha_j}|a_j|_0 .
\]
The upper bound \(f\leq0\) implies \(\sup|a_j|_0\leq1\), so
\eqref{eq:p081-curvebase-unitball} and
\eqref{eq:p081-curvebase-coefficients} apply to every \(a_j\).
Put \(D=mC\). Truncate each coefficient \(c_{j,\nu}(t)\) through degree
\(D\), and let \(b_j\) denote the corresponding finite sum of monomials.
At every \(x\in\mathcal F\), the ultrametric inequality gives
\begin{equation}\label{eq:p081-curvebase-tail}
 \mathrm e^{\alpha_j}|a_j-b_j|_0(x)
 \leq\mathrm e^{\alpha_j-D-1}
 <\mathrm e^{-D}\leq\mathrm e^{-mM}\leq\mathrm e^{mf(x)}.
\end{equation}
There are only finitely many monomials, each of norm at most one.
The scalar coefficient tails belong to \(t^{D+1}R\), which proves this
bound without any no-cancellation assumption at \(x\).

For finitely many pairs \(a_j,b_j\) in a non-Archimedean normed line,
if \(w_j|a_j-b_j|<B=\max_jw_j|a_j|>0\) for every \(j\), then
\(\max_jw_j|b_j|=B\): all terms are at most \(B\), and an index
attaining \(B\) cannot change norm under a strictly smaller error.
Applied in the one-dimensional completed-residue-field fibre of
\(L^{\otimes m}\), this yields
\begin{equation}\label{eq:p081-curvebase-exact-tail}
             f=\frac1m\max_j\{\log|b_j|_0+\alpha_j\}
                         \quad\hbox{on all of }\mathcal F.
\end{equation}
The lower bound on \(f\) makes \(B>0\). Points with proper support and
cancellation among retained terms are included. The coefficients of the
original sections can be arbitrary Laurent series; only their finitely
many tails have been discarded below a strict margin.

\emph{Homogenization and negative padding remove all base loci.}
Write \(b_j=\sum_{n=0}^D a_{j,n}t^n\), where each \(a_{j,n}\) is a
\(k\)-linear combination of global monomials \(q^\nu\). In
\(H^0(S,B_C^{\otimes m})\), take
\[
 B_j=\sum_{n=0}^D a_{j,n}\pi^*(T^nU^{D-n}),\qquad
 Q_{\nu,0}=q^\nu\pi^*U^D,\qquad
 Q_{\nu,\infty}=q^\nu\pi^*T^D .
\]
Give \(B_j\) weight \(\alpha_j\), and give every padding section
weight \(-mC\). The padding system alone generates \(B_C^{\otimes m}\):
at every scheme point some \(q^\nu\) is nonzero and at least one of
\(T,U\) is nonzero. In particular this removes central-fibre base loci,
all other finite-fibre base loci, and those at infinity. The weighted
maximum defines \(\Phi\in H_{\mathbf R}(B_C)\).

At \(\eta\), \(|U|_0=1\) and \(|T|_0=\mathrm e^{-1}\). In the
norm-one frame \(U^D\) for the base factor, \(B_j\) restricts to \(b_j\).
Every weighted padding term has norm at most
\(\mathrm e^{-mC}\leq\mathrm e^{mf}\); hence
\eqref{eq:p081-curvebase-exact-tail} gives \(\Phi|_{\mathcal F}=f\).
Globally each section has trivial norm at most one, while the maximum
of the unweighted padding norms is exactly one. Therefore
\(-C\leq\Phi\leq1/m\leq1\). Its normalized weights are \(-C\) and
\(\alpha_j/m\in[0,1)\). All weights remain rational if the input
weights are rational.

Finally, \(B_C\) is ample: tensoring an ample line bundle with a
globally generated one preserves ampleness. For completeness, choose
a power of the ample bundle which is very ample; its embedding,
together with the morphism defined by the corresponding power of the
globally generated bundle, is a closed immersion into a product of
projective spaces. The Segre embedding shows that their tensor product
power is very ample. Apply this with \(\pi^*\mathcal O(C)\).
\end{proof}

\begin{proof}[Proof of \cref{thm:p081-curvebase-ep-transfer-candidate}]
\emph{A scaling collar of the entire fibre.}
On \(\mathcal C=\{0<|t|<1\}\subset S^{\mathrm{an}}\), put
\(a(y)=-\log|t(y)|>0\). The base point is
\(a(y)\operatorname{ord}_{t=0}\): a nonzero polynomial over trivial
\(k\) has a unique lowest nonzero exponent, whose norm strictly
dominates the other terms when \(0<|t|<1\).
Continuous positive scaling of semivaluations
\cite[\S1.1.4]{BJ22} gives inverse homeomorphisms
\begin{equation}\label{eq:p081-curvebase-collar}
 \mathcal F\times\mathbf R_{>0}\longrightarrow\mathcal C,\quad
 (x,a)\longmapsto ax,\qquad
 y\longmapsto(a(y)^{-1}y,a(y)).
\end{equation}
The fibre identification uses the completed residue field
\(\mathscr H(\eta)=K\)
\cite[Definition/Exercise~4.1.7.1 and \S5.1.1]{Temkin2011}.
Explicitly, affine-chart seminorms over \(\eta\) extend from \(k(t)\)
to its completion \(K\), and give the charts of \(Y_K^{\mathrm{an}}\);
restriction is inverse. The natural map is continuous, and is a
homeomorphism onto the subspace fibre: its source is compact by
projectivity and its target Hausdorff. This is the full analytic fibre,
not the image of a Gauss section. Nothing requires that the family
\(\pi\) be constant.

\emph{The uniform radial estimate is the missing usc ingredient.}
For a family \(f_i\in[-M,0]\), choose \(\Phi_i\) as above using one
\(C\). The trivial metric is homogeneous under scaling:
\(\log|s|_0(ax)=a\log|s|_0(x)\), since the centre is unchanged
\cite[Remark~1.24 and \S1.1.4]{BJ22}. Each normalized weight has
absolute value at most \(C\). Comparison of the finite maxima gives
\begin{equation}\label{eq:p081-curvebase-radial}
 |\Phi_i(ax)-a\Phi_i(x)|\leq C|a-1|,\qquad
 |\Phi_i(ax)-\Phi_i(x)|\leq2C|a-1| \quad(x\in\mathcal F).
\end{equation}
The second estimate uses \(|f_i|\leq M\leq C\). Terms equal to
\(-\infty\) cause no difficulty, since each system has a finite term.
The bound is uniform in the index and the FS presentation.

Put \(V=\sup_i\Phi_i\) and \(g=(V|_{\mathcal F})^{*\mathcal F}\).
Taking suprema in \eqref{eq:p081-curvebase-radial} gives the same
second estimate for \(V\). At \(x\in\mathcal F\), choose a fibre
neighbourhood \(W\) on which \(V|_{\mathcal F}\leq g(x)+\epsilon\).
For \(|a-1|<\epsilon/(2C)\), the collar neighbourhood in
\eqref{eq:p081-curvebase-collar} then has \(V\leq g(x)+2\epsilon\).
This proves \(V^{*S}(x)\leq g(x)\). The reverse inequality follows
by restricting ambient neighbourhoods to the fibre, proving
\eqref{eq:p081-curvebase-usc}. These are neighbourhood arguments,
valid without metrizability or countable cofinal subsets.

Restriction to a closed fibre alone would give only the reverse
inequality. For example, the continuous collar functions
\(\min\{1,n|a-1|\}\) vanish on the fibre, but their supremum has
ambient usc regularization equal to one there. This is a topological
warning, not an asserted FS/psh counterexample in a fixed class.
Our uniform weight estimate rules out that nearby-fibre overshoot.

\emph{Apply ambient EP, then restrict actual FS approximants.}
The \(\Phi_i\) are psh in \(c_1(B_C)\) and globally between \(-C\)
and one. Ambient EP therefore makes
\(\Psi=V^{*S}\) psh in that class. Since \(B_C\) is ample,
\cite[Theorem~4.15(iii) and Example~3.9]{BJ22} produce a decreasing
net in \(H^{\mathrm{dom}}(B_C)\subset H_{\mathbf Q}(B_C)\)
with pointwise limit \(\Psi\). These are FS functions on the actual
fixed bundle \(B_C\), not merely on varying numerical classes.
Each globally generating section system stays generating on \(X\).
Trivializing its base factor by \(U^{mC}\), of norm one on
\(\mathcal F\), gives an FS potential on \(L\), with the reference
metric specified above. Thus the restricted net decreases to
\(\Psi|_{\mathcal F}=g\), and \(g\) belongs to
\(\operatorname{PSH}_{\mathrm{FS}}(L)\).
This proves the envelope assertion for bounded FS families, without
assuming a general singular-psh restriction theorem over \(K\).

\emph{Closure facts which do not assume EP.}
The target class is stable under addition of constants, finite maxima,
and decreasing limits not identically \(-\infty\). For finite maxima,
use products of approximating nets and concatenate section systems
after passing to common tensor powers
\cite[Proposition~5.4]{BE2021}. For decreasing limits use the compact
diagonal lemma \cite[Lemma~7.4 and Corollary~7.5]{BE2021}.
To recall its mechanism, if \(v_i\downarrow v\) and
\(v_{i,j}\downarrow v_i\) are continuous FS approximants, the functions
\(v_{i,j}+1/n\), ordered by reverse pointwise inequality, form a
directed set: choose a common later \(i\) and then use compactness
to get below two prescribed continuous upper bounds with strict
positive margins. Their infimum is \(v\).
Constants are FS relative to the chosen reference, since the
\(q_\ell\) generate \(A\) and \(\max_\ell|q_\ell|_0=1\).
All target functions are usc as decreasing limits of continuous ones.

\emph{Bounded psh families by a product of directed sets.}
Let \(-M\leq u_i\leq0\) be a nonempty family in the target class.
Choose FS nets \(u_{i,j}\downarrow u_i\).
They are already at least \(-M\). Each has a cofinal tail bounded
above by one: the open sets \(\{u_{i,j}<1\}\) cover compact
\(\mathcal F\); choose a finite subcover and a common later index.
Restrict the directed set \(J_i\) to that tail.
The coordinatewise ordered product \(J=\prod_iJ_i\) is directed.
For \(\boldsymbol j=(j_i)\), set
\[
             W_{\boldsymbol j}=(\sup_i u_{i,j_i})^{*\mathcal F}.
\]
The bounded-FS case applies after subtracting one, so each
\(W_{\boldsymbol j}\) is psh. They decrease and remain at least
\(-M\). Their infimum \(W\) is therefore psh by decreasing-net closure.
If \(s=\sup_i u_i\), then \(W\geq s^{*\mathcal F}\).

For the opposite inequality, take any continuous majorant \(v\geq s\)
and \(\epsilon>0\). For each \(i\), the sets
\(\{u_{i,j}<v+\epsilon\}\) cover \(\mathcal F\). Compactness and
directedness give an index \(j_i\) in the selected tail with
\(u_{i,j_i}\leq v+\epsilon\) everywhere. Choose these indices for
all \(i\). Then \(W_{\boldsymbol j}\leq v+\epsilon\), hence \(W\leq v\).
For every bounded function \(s\) on a compact Hausdorff space,
\begin{equation}\label{eq:p081-curvebase-majorants}
       s^*=\inf\{v\in C^0(\mathcal F):v\geq s\}.
\end{equation}
Indeed, continuous majorants dominate \(s^*\). At any fixed point \(x\),
upper semicontinuity supplies a neighbourhood on which
\(s^*<s^*(x)+\delta\). A continuous function equal to zero at \(x\)
and one outside that neighbourhood exists by normality of the compact
Hausdorff space. Adding a sufficiently large nonnegative multiple of
it to \(s^*(x)+\delta\) gives a continuous majorant. Let \(\delta\)
decrease to zero. Equation \eqref{eq:p081-curvebase-majorants} proves
\(W=s^{*\mathcal F}\). The coordinate choices use ordinary
set-theoretic choice, not an interchange of arbitrary infima and suprema.

\emph{Remove the common lower bound.}
For a generally bounded-above family, first subtract a common upper
bound. The clamps \(\max\{u_i,-M\}\) are psh and lie in \([-M,0]\).
The result just proved and the neighbourhood definition of usc
regularization give
\[
 (\sup_i\max\{u_i,-M\})^{*\mathcal F}
          =\max\{(\sup_i u_i)^{*\mathcal F},-M\}
            \in\operatorname{PSH}_{\mathrm{FS}}(L).
\]
As positive integers \(M\) increase these decrease to the required
envelope. It is not identically \(-\infty\), since it majorizes a member
of the nonempty family. Decreasing-net closure finishes the proof.
No individual unbounded psh function was lifted to the ambient space.
\end{proof}

\begin{proof}[Proof of \cref{prop:p081-curvebase-arrange}]
\emph{A normal projective model retaining a power of the polarization.}
Choose \(n>0\) with \(L_0^{\otimes n}\) very ample, embed
\(Y\) in \(\mathbf P^r_{F_0}\) by it, and take the scheme-theoretic
closure \(S_0\subset\mathbf P^r_k\times\mathbf P^1_k\).
This closure is integral and projective, and its generic fibre is \(Y\).
Normalize it, \(\nu:S\to S_0\).
Finite-type schemes over \(k\) are excellent, hence Nagata
\cite[Tag 07QS]{Stacks}; the normalization is finite
\cite[Tag 035S]{Stacks}. Consequently \(S\) is normal, integral and
projective. Integral closure commutes with localization (a monic
equation can have its denominators cleared), so its generic fibre is
the normalization of \(Y\), which is \(Y\) since \(Y\) is smooth.
The resulting dominant map to \(\mathbf P^1\) is flat by the
torsion-free DVR argument already given.

The line bundle
\[
          B_0=\nu^*\bigl((\mathcal O_{\mathbf P^r}(1)
                     \boxtimes\mathcal O_{\mathbf P^1}(1))|_{S_0}\bigr)
\]
is ample by the Segre embedding and finite pullback of ampleness
\cite[Tag 0B5V]{Stacks}. Its generic restriction is
\(L_0^{\otimes n}\), after choosing a frame for the base factor.
Choose \(r_0>0\) such that \(B=B_0^{\otimes r_0}\) is very ample
over \(k\). Its global sections give a closed immersion
\(j_0:S\hookrightarrow\mathbf P^r_k\), with \(r\) now rechosen.
The graph map
\[
        j=(j_0,\pi):S\hookrightarrow
                          \mathbf P^r_k\times\mathbf P^1_k
\]
is a closed immersion, and \(j^*\mathcal O(1,0)=B\).

\emph{Relative Serre vanishing gives all the required multiplication.}
Let \(\mathcal I\) be the coherent ideal of this graph, and let
\(p:\mathbf P^r\times\mathbf P^1\to\mathbf P^1\) be projection.
On each of the two standard affine base charts, apply Serre vanishing
to \(\mathcal I\) twisted by \(\mathcal O(n,0)\), using
\cite[Tag 0B5T(4)]{Stacks}. This bundle is ample on the projective
space over that affine chart. Take one \(n_0\) working for both charts.
The exact ideal sequence and its long exact cohomology sequence show
that for all \(n\geq n_0\)
\begin{equation}\label{eq:p081-curvebase-poly-surjection}
 \operatorname{Sym}^n(k^{r+1})\otimes_k\mathcal O_{\mathbf P^1}
                     \longrightarrow\pi_*B^{\otimes n}
                 \quad\text{is surjective}.
\end{equation}
More explicitly, on an affine chart the source is the degree-\(n\)
homogeneous polynomials, and the vanishing of
\(H^1(\mathcal I(n,0))\) makes its map onto
\(H^0(S_{\mathrm{chart}},B^{\otimes n})\) surjective.
These are the local module descriptions of the displayed sheaf map.
No specialization base-change theorem and no unproved assertion about
an arbitrary Veronese algebra is being used.

Choose \(d\geq\max\{1,n_0\}\), and set \(A=B^{\otimes d}\).
It is very ample, hence ample and globally generated, and
\eqref{eq:p081-curvebase-poly-surjection} for \(n=d\) makes
\(E_1=\pi_*A\) globally generated. Every degree-\(md\) monomial in
the \(r+1\) variables is a product of \(m\) degree-\(d\) monomials.
Thus the multiplication of polynomial systems in degree \(d\), followed
by their images in \(E_1\), maps onto \(E_m=\pi_*B^{\otimes md}\)
by \eqref{eq:p081-curvebase-poly-surjection}. This factors through
\(\operatorname{Sym}^m E_1\), proving its required surjectivity for
every \(m\geq1\). Take \(N=nr_0d\); then \(A|_Y\simeq L_0^{\otimes N}\).
Finite systems generating both \(E_1\) and \(A\) can be chosen by
combining their global generating systems.

\emph{Return from the tensor power to the original line bundle.}
Apply \cref{thm:p081-curvebase-ep-transfer-candidate} to \(A|_X\).
Its model reference divided by \(N\) is an FS metric on \(L_{0,K}\):
the generators of \(A\) give its FS presentation in tensor degree \(N\).
Multiplication by \(N\) identifies FS metrics on \(L_{0,K}\) with FS
metrics on \(L_{0,K}^{\otimes N}\)
\cite[Proposition~5.4(iii)]{BE2021}; this also follows directly by
raising a section presentation to common tensor powers. It therefore
identifies their decreasing-FS classes. Supremum and usc
regularization commute with multiplication by a positive constant,
so the envelope assertion descends from the tensor power.
Changing a continuous reference subtracts the same continuous
function from all potentials and commutes with usc regularization;
thus this is a statement about metrics, not only this chosen reference.

This construction uses a projective model of an already
\(k(t)\)-defined polarized variety. It never replaces arbitrary
equations over \(k((t))\) by polynomial truncations: that would change
the variety and is not justified by the section-tail argument.
\end{proof}
